% Generated draft: requires mathematical and source review.
\documentclass[11pt]{article}
\usepackage[margin=25mm]{geometry}
\usepackage{fontspec}
\setmainfont{Latin Modern Roman}
\usepackage{amsmath,amssymb,mathtools,unicode-math}
\setmathfont{Latin Modern Math}
\usepackage{enumitem,needspace,xurl,hyperref}
\setlength{\parindent}{0pt}
\setlength{\parskip}{5pt}
\setlength{\emergencystretch}{3em}
\newcommand{\N}{\mathbb{N}}
\newcommand{\Z}{\mathbb{Z}}
\newcommand{\Q}{\mathbb{Q}}
\newcommand{\R}{\mathbb{R}}
\newcommand{\C}{\mathbb{C}}
\newcommand{\F}{\mathbb{F}}
\newcommand{\A}{\mathbb{A}}
\newcommand{\PP}{\mathbb{P}}
\newcommand{\E}{\mathbb{E}}
\newcommand{\eps}{\varepsilon}
\newcommand{\dd}{\,\mathrm{d}}
\begin{document}
\section*{UnsolvedMath ID: 30002772}
\textbf{Source problem code:} catalogue entry
\subsection{Definitions and mathematical statement}
Let \(M\) be a complex manifold and let \(\Omega\subseteq M\) be a domain, meaning a nonempty connected open subset. A plurisubharmonic function is an upper semicontinuous function \(u:\Omega\to[-\infty,\infty)\) whose pullback along every holomorphic disc is subharmonic or identically \(-\infty\). It is bounded above if \(\sup_{\Omega}u<\infty\).

The supplied record describes \(c(\Omega)\) as the set where the Levi form of every bounded-above plurisubharmonic function degenerates. This requires a regularity qualification: an arbitrary plurisubharmonic function need not have a pointwise classical Levi form. For each \(r\in\{2,\infty\}\), a precise candidate interpretation is obtained by setting
\[
\mathcal P_r(\Omega)=\{u\in C^r(\Omega,\R):u\text{ is plurisubharmonic and }\sup_{\Omega}u<\infty\}.
\]
For \(u\in\mathcal P_r(\Omega)\), its Levi form in local holomorphic coordinates is
\[
\mathcal L_u(p;v)=\sum_{j,k=1}^{n}\frac{\partial^2u}{\partial z_j\partial\overline z_k}(p)\,v_j\overline v_k,
\qquad v\in T_p^{1,0}M.
\]
Since this form is positive semidefinite, degeneracy means that it has a nonzero null vector. Thus the corresponding candidate core is
\[
c_r(\Omega)=\left\{p\in\Omega:\ \forall u\in\mathcal P_r(\Omega)\ \exists v\in T_p^{1,0}M\setminus\{0\},\ \mathcal L_u(p;v)=0\right\}.
\]
The null vector may depend on \(u\). The record does not specify whether \(c(\Omega)\) means \(c_2(\Omega)\), \(c_\infty(\Omega)\), or a definition using nonsmooth functions. No equality of these interpretations is asserted here.

A set \(E\subseteq\Omega\) is locally pluripolar if every \(p\in\Omega\) has a connected open neighborhood \(U\subseteq\Omega\) and a plurisubharmonic function \(\varphi\) on \(U\), not identically \(-\infty\), such that
\[
E\cap U\subseteq\{\varphi=-\infty\}.
\]
It is globally pluripolar in \(\Omega\) if a single such function on \(\Omega\) contains all of \(E\) in its \(-\infty\) locus. The source does not specify which meaning of pluripolarity is intended. It also does not specify the convention for pseudoconvexity on an arbitrary complex manifold; in particular, no boundary regularity or plurisubharmonic exhaustion hypothesis is supplied.

\textbf{Problem, as supplied.} Understand how geometric properties of \(\Omega\) are related to properties of its core \(c(\Omega)\). In particular, determine what relationship holds between pseudoconvexity of \(\Omega\) and pluripolarity of \(c(\Omega)\). This is a research question, not a specified implication or equivalence. A definitive theorem statement requires clarification of the core's function class, the meanings of pseudoconvexity and pluripolarity, and any necessary restrictions on the ambient manifold or domain.
\subsection{Short English statement}
The question asks how the geometry of a complex domain controls the set where every admissible bounded-above plurisubharmonic function fails to be strictly plurisubharmonic. It particularly asks about pseudoconvexity and pluripolarity of this core. The supplied formulation leaves essential conventions unspecified, so the displayed core definitions are conditional interpretations requiring review.
\subsection{Sources}
\begin{itemize}
\item Local UnsolvedMath registry, permanent ID \(30002772\), title \emph{Pseudoconvex Domains and Their Cores}; exact source problem code \texttt{catalogue entry}. The original recorded question is to understand how properties of \(\Omega\) relate to properties of \(c(\Omega)\). The record gives proposed year 2015 and names no proposer.
\item Supplied report citation: \emph{Geometric Methods of Complex Analysis} (2015), Oberwolfach Reports, pp. 235--283, DOI \texttt{10.4171/owr/2015/4}. Supplied literature-source label: ``Merker--Porten, Oberwolfach Report.'' \url{https://doi.org/10.4171/owr/2015/4}. This label is preserved as attribution of the reference; it does not establish who proposed the problem.
\item Supplied literature source: ``MFO primary report.'' \url{https://publications.mfo.de/bitstream/handle/mfo/3451/OWR_2015_04.pdf?isAllowed=y&sequence=1}.
\item Supplied literature assessment, checked on \texttt{2026-08-21}: status ``partially solved.'' The assessment reports that the core is proved \(1\)-pseudoconcave and that specific relations between domain geometry and the core are established, while the general relationship remains open. This is supplied dated curation, not independently verified here; the hypotheses and precise versions of those reported results are not supplied.
\item The registry labels the statement \texttt{corrected\_verified} and reports that the official report passage was checked directly and its question, definitions, and notation restored without a stronger claim. This is a supplied source-verification note, not a claim that this draft independently inspected the report. The background identifies Oberwolfach Reports as records of mathematical workshops and their open-problem discussions and expressly requests independent verification of time-sensitive status and cited claims.
\end{itemize}
\end{document}
